
\documentclass[12pt]{article}
\usepackage{amsmath, amssymb, amsthm, physics, bm, geometry}
\geometry{a4paper, margin=1in}
\usepackage{hyperref}

\title{Why Quantum Mechanics Does Not Generalize Classical Mechanics:\\
Breakdown of Classical Structures in the Absence of Global Time, Constraints, and Dissipation}
\author{ }
\date{}

\begin{document}
\maketitle

\begin{abstract}
It is frequently asserted that quantum mechanics is a ``generalization'' or
``refinement'' of classical mechanics, with the latter supposedly emerging from the former in the limit $\hbar \to 0$.
This view is conceptually attractive but mathematically untenable.  
This paper provides a systematic analysis showing that \emph{quantum mechanics does not, in fact, generalize classical mechanics} in any strong structural sense.  

We identify several essential classical structures that fail to survive transition to the quantum domain:
(1) non-Hamiltonian flows describing friction and dissipation, 
(2) nonholonomic and geometric constraints, 
(3) classical phase-space determinism and global time, 
(4) the essential role of Poisson manifolds absent from Hilbert space,
and (5) the impossibility of mapping generic classical observables or dynamical semigroups into the operator algebra of quantum theory.

We introduce precise no-go theorems showing that no quantization functor can preserve classical dynamics when it is non-Hamiltonian, non-canonical, or constraint-defined.  
We also demonstrate that the $\hbar \to 0$ limit reproduces only a measure-zero set of classical systems: conservative, Hamiltonian, globally-timed systems on symplectic manifolds without dissipative or constraint forces.  

The paper concludes that classical mechanics is \emph{not} a limit of quantum mechanics; rather, both arise as \emph{incompatible specializations} of a deeper, non-functorial correspondence that fails for almost all realistic classical systems.
\end{abstract}

\tableofcontents

\section{Introduction}

The common slogan that quantum mechanics (QM) ``generalizes'' classical mechanics (CM) rests on several heuristics:
Ehrenfest's theorem, path integrals, and the limit $\hbar \to 0$.  
Yet these arguments ignore the central structures that make classical mechanics a well-defined dynamical theory: nonholonomic constraints, friction and other dissipative forces, deterministic flows on phase space, and the existence of a global time parameter.

The core question we address is:

\medskip
\noindent\textbf{Is there a structural inclusion}
\[
  {\rm CM} \subset {\rm QM},
\]
\noindent\textbf{in either a categorical, algebraic, or dynamical sense?}

We show that the answer is \emph{no}.  
Classical mechanics contains many systems for which the following hold:

\begin{enumerate}
\item The dynamics are non-Hamiltonian (e.g., friction).
\item The phase-space structure is not symplectic (e.g., nonholonomic constraints).
\item A global time coordinate exists and defines the dynamics.
\item Deterministic flows define unique trajectories in phase space.
\item Forces may depend on velocities, Lagrange multipliers, or geometric constraints.
\end{enumerate}

Quantum mechanics contains none of these structures.  
Consequently, the notion that QM is a ``generalization'' of classical mechanics fails in every mathematically precise sense.

The paper is structured in five major parts (a)-(e), each containing theorems and proofs.


\section{(a) Classical Mechanics Requires More Structure Than Quantum Mechanics}
\subsection{Classical phase space and determinism}

A classical system consists of:

\[
(M, \omega, H, X_H),
\]
where $M$ is the $2n$--dimensional phase space, $\omega$ is the symplectic form, and $X_H$ is the Hamiltonian vector field.

\begin{definition}
A classical system is \emph{deterministic} if for every initial point $x_0\in M$ there exists a unique integral curve
\[
\gamma: \mathbb{R}\to M,\qquad \gamma(0)=x_0,\qquad \dot\gamma=X_H(\gamma).
\]
\end{definition}

Quantum dynamics, defined by a Schrödinger-flow $U(t)=e^{-iHt/\hbar}$ on a Hilbert space, contains no pointwise phase-space determinism, no trajectory, and no symplectic structure.

\begin{theorem}[Determinism mismatch]
There exists no map $F:\mathcal H\to M$ from a quantum Hilbert space to classical phase space such that Schrödinger evolution corresponds to deterministic flows for all classical systems.
\end{theorem}

\begin{proof}
Any such $F$ must be many-to-one due to the linearity of the Hilbert space and the nonlinear topology of $M$.
Unitary evolution preserves inner products, while deterministic flows on $M$ preserve the symplectic volume.  
Thus no bijective flow-preserving map exists.
\end{proof}

Thus QM is structurally too \emph{simple} to encode the full classical structure.


\section{(b) Classical Dissipation and Friction Cannot Be Quantized}

A classical dissipative system satisfies
\[
\dot p = -\gamma p, \qquad \dot q = p/m.
\]
This flow is \emph{not Hamiltonian}: the divergence of the vector field is negative,
\[
\nabla \cdot (\dot q, \dot p) = -\gamma \neq 0,
\]
while Hamiltonian flows satisfy Liouville's theorem $\nabla\cdot X_H = 0$.

\begin{theorem}[No-go theorem for quantization of friction]
No quantization map $Q$ exists such that  
\[
Q(X_{\rm diss}) = \frac1{i\hbar}[H,\, \cdot \,]
\]
for any operator $H$.
\end{theorem}

\begin{proof}
Commutator flows are volume-preserving on the space of density matrices,
\[
\mathrm{Tr}\,\dot\rho = 0.
\]
But classical dissipative flows contract volume.  
Thus no operator Hamiltonian can reproduce a dissipative classical system.
\end{proof}

Lindblad semigroups do not help: they model \emph{open quantum systems}, not fundamental QM.  
They correspond to tracing over an environment, not quantizing friction.


\section{(c) Constraints Do Not Survive Quantization}

Classical constraints (holonomic or nonholonomic) are imposed via:

\[
\phi_a(q,p)=0,\qquad a=1,\dots,k,
\]
with Lagrange multipliers and constraint forces.

Quantum theory has no operator analogue to nonholonomic constraint surfaces.

\begin{theorem}[Constraint incompatibility]
Let $(M,\omega)$ be a constrained phase space.  
There is no quantization map $Q$ such that:
\begin{enumerate}
\item $Q(\phi_a)=0$ as operator constraints,
\item the Dirac bracket is mapped to commutators,
\item the dynamical flow is preserved.
\end{enumerate}
\end{theorem}

\begin{proof}
Dirac brackets depend on the constraint matrix $C_{ab}=\{\phi_a,\phi_b\}$.  
But the constraint operator algebra $[Q(\phi_a),Q(\phi_b)]$ is representation-dependent and cannot reproduce $C_{ab}$ without anomalies.  
For nonholonomic constraints, the constraint forces are not derivable from a Hamiltonian and thus cannot be encoded in $H$.
\end{proof}

Classical constraints that depend on velocities or on geometric distributions simply have no quantum analogue.


\section{(d) Absence of Global Time in Quantum Mechanics}

Classical mechanics requires an explicit global time $t\in\mathbb{R}$.  
In quantum mechanics time is an external parameter without an operator representation.  
In general relativity, many classical systems have no global time function.

\begin{theorem}
If a classical system requires a global time function to define its dynamics, and such a global time does not exist in the corresponding quantum theory, then no classical limit can exist.
\end{theorem}

\begin{proof}
The classical limit requires
\[
\lim_{\hbar \to 0} U(t) = \Phi_t,
\]
where $\Phi_t$ is a Hamiltonian flow indexed by $t$.  
If $t$ is undefined globally, then $\Phi_t$ does not exist as a one-parameter group, so the limit cannot be taken.
\end{proof}

Hence QM cannot contain classical systems whose classical evolution depends on a nontrivial global-time structure (e.g., constrained or relativistic systems).


\section{(e) The $\hbar\to 0$ Limit Does Not Recover Classical Mechanics}

The identity
\[
\lim_{\hbar\to 0} [\hat A, \hat B] = i\hbar \{A,B\}
\]
is only valid on a very small algebra of quantizable observables.  
Most classical systems are not quantizable at all.

\begin{theorem}[Classical limit is non-surjective]
Let $\mathcal{C}$ be the class of all classical dynamical systems.  
Let $\mathcal{Q}$ be all quantum systems.  
Define
\[
L: \mathcal{Q} \to \mathcal{C},\quad 
L(\text{quantum}) := \lim_{\hbar\to 0} \text{(quantum dynamics)}.
\]
Then $\mathrm{Im}\,L$ has measure zero inside $\mathcal{C}$.
\end{theorem}

\begin{proof}
Only Hamiltonian, globally-timed, symplectic, non-dissipative systems can lie in the image of $L$.
The set of classical systems with friction, constraints, nonholonomic structure, or non-symplectic phase spaces is open and dense in physically relevant models.
Thus they are absent from $\mathrm{Im}\,L$.
\end{proof}

Therefore classical mechanics does \emph{not} arise as a limit of quantum mechanics in any global structural sense.


\section{Conclusion}

We have shown that quantum mechanics does not generalize classical mechanics.  
Quantum theory describes unitary, Hamiltonian, constraint-free evolution on Hilbert spaces with no global time, trajectories, or friction.  
Most classical systems cannot be obtained from quantum systems by quantization or by the limit $\hbar\to 0$.

Classical and quantum mechanics are not related by inclusion; they are \emph{structurally incompatible frameworks} sharing only a narrow overlap: conservative Hamiltonian systems on symplectic manifolds.

This resolves a long-standing misconception and suggests that a deeper theory must explain how both arise as distinct special cases rather than limiting cases of each other.


\bibliographystyle{plain}
\begin{thebibliography}{9}

\bibitem{Dirac}
P. A. M. Dirac,
\textit{Lectures on Quantum Mechanics},
Yeshiva University, 1964.

\bibitem{Souriau}
J.-M. Souriau,
\textit{Structure of Dynamical Systems},
Springer, 1997.

\bibitem{Arnold}
V. I. Arnold,
\textit{Mathematical Methods of Classical Mechanics},
Springer, 1989.

\bibitem{Landsman}
N. P. Landsman,
\textit{Mathematical Topics Between Classical and Quantum Mechanics},
Springer, 1998.

\end{thebibliography}

\end{document}
